\chapter*{Перечень обозначений и сокращений}
\addcontentsline{toc}{chapter}{Перечень обозначений и сокращений}

\begingroup
\setlength{\parindent}{0pt}
\setlength{\parskip}{0.45\baselineskip}

\textbf{RAG} Retrieval-Augmented Generation — подход к генерации текста, при котором языковая модель использует контекст, предварительно найденный во внешней базе знаний.

\textbf{LLM} Large Language Model — большая языковая модель, предназначенная для обработки и генерации текста на естественном языке.

\textbf{NLP} Natural Language Processing — обработка естественного языка.

\textbf{API} Application Programming Interface — программный интерфейс взаимодействия компонентов.

\textbf{Embedding} Числовое векторное представление текста, отражающее его смысловые признаки в многомерном пространстве.

\textbf{Vector database} Векторная база данных, предназначенная для хранения embeddings и поиска ближайших векторов.

\textbf{Qdrant} Векторная база данных, используемая в проекте для хранения embeddings чанков и связанных metadata.

\textbf{Chunk} Фрагмент документа, используемый как минимальная единица поиска и передачи контекста в LLM.

\textbf{Retrieval} Этап поиска релевантных фрагментов по пользовательскому вопросу.

\textbf{Dense retrieval} Семантический поиск, при котором вопрос и документы сравниваются по близости их векторных представлений.

\textbf{Prompt} Текстовая инструкция и контекст, передаваемые языковой модели для генерации ответа.

\textbf{JSONL} Формат хранения структурированных данных, в котором каждая строка является отдельным JSON-объектом.

\textbf{Hit@K} Метрика поиска, показывающая долю вопросов, для которых хотя бы один релевантный фрагмент найден среди первых $K$ результатов.

\textbf{MRR@K} Mean Reciprocal Rank — среднее обратное значение ранга первого релевантного результата среди первых $K$ результатов.

\textbf{Recall@K} Метрика полноты поиска, показывающая долю релевантных фрагментов, найденных среди первых $K$ результатов.

\textbf{MVP} Minimum Viable Product — минимальная версия системы, достаточная для проверки базового сценария работы.
\endgroup

%%% Local Variables:
%%% TeX-engine: xetex
%%% End:
